\section{Related work}
\label{sec:related}

\paragraph{Reward hacking and proxy gaming.} That optimising a proxy past a point degrades the
true objective is established theoretically \citep{skalse2022defining}, empirically for learned
reward models \citep{gao2023scaling}, and across hand-specified rewards
\citep{amodei2016concrete,pan2022effects}. The consistent framing is volumetric --- how much
proxy score is gained, how much true objective is lost, and where the divergence begins. Our
contribution is narrower and, we think, under-reported: what happens to the \emph{composition}
of the degradation when a mitigation closes one route to it while the proxy's exploitable gap
stays open.

\paragraph{Length and sycophancy biases in LLM judges.} LLM raters systematically prefer longer
outputs \citep{singhal2024long,dubois2024lengthcontrolled,zheng2023judging} and agreeable ones
\citep{sharma2024sycophancy}, and prefer their own family's text \citep{panickssery2024selfpreference}.
Our flattery result is a clinical instance of the sycophancy bias, with the twist that the
instruments being gamed are published clinical questionnaires rather than generic preference
ratings --- and our verbosity analysis (\S\ref{sec:aggregate}) shows the length bias operating
through a moving denominator rather than through the score alone.

\paragraph{Trajectory-level and process-level reward.} Scoring a decision by the continuation it
leads to, rather than in isolation, underlies look-ahead in preference-tree optimisation
\citep{baruch2025pto}, MCTS-guided preference construction \citep{xie2024mctsdpo,yu2023promptmcts},
efficient planning over simulated conversation futures \citep{chen2025broaden}, and, one level
down, process supervision over reasoning steps \citep{lightman2024letsverify}. These methods are
usually motivated as variance or credit-assignment improvements. We evaluate look-ahead instead
as a \emph{reward-hacking mitigation} --- a use it is often informally credited with and, to our
knowledge, has not been tested as.

\paragraph{NLP for Motivational Interviewing.} MI comes with an official treatment-integrity
coding system and published competency thresholds \citep{moyers2016miti}, and counselling quality
is separable from language alone \citep{perezrosas2019goodcounselor}. LLM-simulated patients are
used to train human counsellors \citep{steenstra2025scaffolding}, and \citet{yosef2024assessing}
validate the questionnaire-filling LLM evaluators our reward is built from. We use MI fidelity
coding as the \emph{outcome} of an RL study rather than as the modelling target: the coded
behaviour inventory is what lets a substitution be distinguished from a reduction at all.

\paragraph{Multi-judge evaluation.} Agreement between LLM judges is usually treated as
inter-rater reliability \citep{zheng2023judging,koo2016icc}. Ours cannot be: the primary grader
\emph{was} the training signal, so the second judge is a train/test split on the reward, not a
second rater --- their raw scores are never pooled, and their disagreement over the aggregate in
\S\ref{sec:aggregate} is itself a finding about aggregates over heterogeneous failures, not a
reliability failure to be averaged away.
